\clearpage
\section{Третий промежуточный отчёт}

Третий этап учебной практики был посвящён сериализации объектов, организации параллельных вычислений и экспериментальному определению параметров многопоточного алгоритма. Работы выполнялись последовательно: сначала был изучен способ сохранения структурированных данных в формате JSON, затем реализован параллельный метод трапеций, после чего проведено измерение его точности и производительности. Во всех заданиях использовались модульные тесты и автоматическая проверка проекта.

\begin{table}[h]
    \centering
    \caption{Состав работ третьей недели}
    \small
    \begin{tabular}{>{\centering\arraybackslash}p{1.1cm} >{\RaggedRight\arraybackslash}p{8.4cm} >{\centering\arraybackslash}p{2.1cm} >{\centering\arraybackslash}p{2.2cm}}
        \toprule
        \textbf{№} & \textbf{Задание} & \textbf{Тестов} & \textbf{Результат} \\
        \midrule
        13 & Сериализация и десериализация объектов в JSON & 14 & успешно \\
        14 & Знакомство с потоками & 17 & успешно \\
        15 & Определение оптимальных параметров вычисления интеграла & 12 & успешно \\
        \bottomrule
    \end{tabular}
\end{table}

\subsection{Сериализация и десериализация объектов}

В задании были определены классы \codebreak{Student} и \codebreak{Subject}. Объект студента содержит имя, фамилию, дату рождения и список предметов с оценками. Для преобразования объектов применён стандартный механизм \codebreak{System.Text.Json} \cite{systemtextjson}. Работа была выполнена в несколько этапов. Сначала создана модель данных, затем разработан сервис сериализации, после этого добавлены сохранение и загрузка JSON-файла и, наконец, написаны тесты для корректных и ошибочных входных данных.

Параметры сериализатора настроены в одном месте. Свойства со значением \codebreak{null} не включаются в итоговый JSON, поэтому файл не содержит лишних полей. Для даты рождения создан пользовательский конвертер, использующий формат \codebreak{dd.MM.yyyy}. Такой подход обеспечивает одинаковое представление даты независимо от региональных параметров операционной системы. При чтении конвертер принимает только ожидаемый формат и сообщает об ошибке, если строка не является корректной датой.

После десериализации выполняется проверка объекта. Имя и фамилия не должны быть пустыми, дата рождения не может находиться в будущем, название предмета должно быть задано, а оценка должна входить в допустимый диапазон от 1 до 5. Проверка после преобразования необходима потому, что синтаксически правильный JSON ещё не гарантирует предметную корректность данных. Например, значение оценки 20 может быть успешно прочитано как целое число, но не может считаться допустимой оценкой студента.

При выполнении задания возникло несколько затруднений. Первое было связано с тем, что стандартный кодировщик JSON может заменять кириллические символы последовательностями Unicode. Это не нарушает корректность файла, но снижает его читаемость. Проблема была решена настройкой кодировщика, сохраняющего текст в понятном виде. Второе затруднение состояло в разграничении ошибок формата и ошибок предметной области. Неверная дата должна приводить к ошибке преобразования, а пустое имя или недопустимая оценка --- к ошибке валидации. Для этого ответственность была разделена между конвертером и отдельным методом проверки данных.

Наиболее сложным моментом оказалось обеспечение симметрии операций: объект после последовательности «сериализация --- десериализация» должен сохранять все значимые данные. Одновременно требовалось игнорировать отсутствующие значения, использовать специальный формат даты и отклонять недопустимые объекты. Корректность решения подтверждена тестами на обычную сериализацию, отсутствие полей со значением \codebreak{null}, формат даты, сохранение и загрузку файла, повреждённый JSON, будущую дату рождения и оценки вне установленного диапазона.

\subsection{Знакомство с потоками}

Во втором задании недели реализовано вычисление определённого интеграла методом трапеций. Для соседних точек $x_i$ и $x_{i+1}$ площадь элементарной трапеции вычисляется по формуле
\[
    S_i = \frac{f(x_i)+f(x_{i+1})}{2}(x_{i+1}-x_i).
\]
Общий результат представляет собой сумму площадей всех трапеций. Благодаря аддитивности интеграла исходный отрезок был разделён на части по количеству потоков, а каждая часть вычислялась независимо.

Для каждой части создавался объект \codebreak{Thread} \cite{threads}. Основной поток ожидал завершения рабочих потоков при помощи примитива \codebreak{Barrier} \cite{barrier}, а затем дополнительно выполнял соединение потоков, чтобы гарантировать освобождение ресурсов. Локальный результат каждого рабочего потока добавлялся к общей сумме атомарно. Поскольку для типа \codebreak{double} отсутствует простая операция атомарного сложения, использован цикл на основе \codebreak{Interlocked.CompareExchange} \cite{interlocked}. Каждый поток сначала вычисляет результат своего отрезка в локальной переменной и обращается к общей переменной только один раз. Это уменьшает количество синхронизаций.

Особое внимание уделено границам разбиения. Части отрезка не должны пересекаться и не должны оставлять пропуски. Последняя трапеция может иметь ширину меньше заданного шага, поэтому её правая граница ограничивается значением $b$. Также предусмотрены проверки входных данных: функция не может быть равна \codebreak{null}, шаг и количество потоков должны быть положительными, а левая граница не должна превышать правую.

Основная трудность состояла в безопасном суммировании результатов. Обычная операция \codebreak{total += localResult} включает чтение, сложение и запись, поэтому два потока могут потерять результат друг друга. Атомарное обновление устранило данную гонку. Дополнительная сложность возникла при обработке исключений внутри рабочих потоков: исключение нельзя просто передать в основной поток. Для этого ошибки собираются в потокобезопасной очереди, а после завершения вычислений основная функция формирует общее исключение.

Наиболее сложным моментом стала согласованная работа барьера, завершения потоков и обработки ошибок. Каждый рабочий поток должен обязательно достигнуть точки синхронизации даже при возникновении исключения, иначе основная функция будет ожидать бесконечно. Поэтому сигнал барьеру размещён в гарантированно выполняемой части кода. Тесты проверяют функции $x$ и $\sin(x)$, разные количества потоков, обратимые симметричные интервалы, нулевую длину отрезка и некорректные параметры.

\subsection{Определение оптимальных параметров алгоритма}

Эксперимент проводился для функции $\sin(x)$ на отрезке $[-100;100]$ при требуемой точности $10^{-4}$. Теоретическое значение интеграла равно нулю, так как функция нечётная, а отрезок симметричен. Для корректного сравнения дополнительно реализован метод \codebreak{SolveSingleThread}, который выполняет тот же метод трапеций без создания объектов \codebreak{Thread}, барьера и атомарных операций.

Время измерялось классом \codebreak{Stopwatch} \cite{stopwatch}. Каждый вариант запускался три раза, после чего вычислялось среднее арифметическое. Исследованы шаги $10^{-1}$, $10^{-2}$, $10^{-3}$, $10^{-4}$, $10^{-5}$ и $10^{-6}$. Все они обеспечили требуемую точность. Абсолютные ошибки находились приблизительно в диапазоне от $2{,}58\cdot10^{-15}$ до $5{,}02\cdot10^{-13}$. Из заданного набора был выбран минимальный размер шага $10^{-6}$. Он создаёт наибольшее количество трапеций и тем самым позволяет объективнее оценить выгоду параллельного вычисления.

Для выбранного шага программно исследованы варианты от 1 до 16 потоков. Число 16 соответствует количеству логических процессоров использованной системы. Результаты каждого запуска сохранены в CSV-файлы, а графики построены библиотекой ScottPlot \cite{scottplot}. Минимальное среднее время $190{,}761$ мс было получено при 15 потоках. При 16 потоках среднее время увеличилось до $194{,}406$ мс, поэтому оптимальным для данной системы признано значение 15.

Однопоточный алгоритм без использования потоков выполнился в среднем за $1617{,}999$ мс. Следовательно, преимущество лучшего многопоточного варианта составило
\[
    \frac{1617{,}999-190{,}761}{1617{,}999}\cdot100\% = 88{,}21\%.
\]
Полученное ускорение значительно превышает требуемые 15 процентов. Наиболее сложным этапом исследования было обеспечить сопоставимость результатов. Однопоточная и многопоточная версии используют одну формулу и одинаковый шаг, замеры повторяются одинаковое число раз, а на график выводятся средние значения. Благодаря этому измеряется влияние организации вычислений, а не различие математических алгоритмов.

\subsection{Описание и интерпретация результатов эксперимента}

На рисунке~\ref{fig:week3-threads} по оси абсцисс указано среднее время выполнения функции \codebreak{Solve}, а по оси ординат --- количество потоков, как требовалось в задании. При переходе от одного потока к двум время уменьшается с $1597{,}835$ до $982{,}655$ мс. Дальнейшее увеличение количества потоков также ускоряет вычисление, но прирост постепенно становится меньше.

\begin{figure}[ht]
    \centering
    \includegraphics[width=0.94\textwidth]{week3-thread-performance.png}
    \caption{Зависимость среднего времени вычисления от количества потоков}
    \label{fig:week3-threads}
\end{figure}

Форма графика объясняется сочетанием параллельного выполнения и накладных расходов. В начале каждый дополнительный поток получает значительную самостоятельную часть отрезка, поэтому вычислительная нагрузка распределяется между ядрами. Однако ускорение не может быть строго пропорционально числу потоков. Время расходуется на создание потоков, работу планировщика, синхронизацию и итоговое суммирование. Кроме того, потоки конкурируют за вычислительные ресурсы и общие уровни кэш-памяти. Поэтому после определённого момента накладные расходы становятся сопоставимы с выигрышем. Небольшое ухудшение результата при 16 потоках является нормальным проявлением этой зависимости и подтверждает необходимость экспериментального, а не формального выбора параметра.

Рисунок~\ref{fig:week3-comparison} показывает непосредственное сравнение реализаций. Высота столбца многопоточного варианта существенно меньше. Разница в $1427{,}238$ мс соответствует ускорению на $88{,}21\%$. Для шага $10^{-6}$ на отрезке длиной 200 выполняется около 200 миллионов вычислительных шагов. При такой нагрузке полезная работа значительно превосходит затраты на запуск и синхронизацию потоков, вследствие чего параллельный вариант оказывается эффективным.

\begin{figure}[ht]
    \centering
    \includegraphics[width=0.88\textwidth]{week3-single-vs-multi.png}
    \caption{Сравнение однопоточной и оптимальной многопоточной реализаций}
    \label{fig:week3-comparison}
\end{figure}

Результаты исследования шага требуют отдельной интерпретации. Обычно уменьшение шага повышает точность метода трапеций, но в проведённом опыте даже шаг $10^{-1}$ дал очень малую ошибку. Причиной является специальный выбор функции и интервала: значения нечётной функции в симметричных точках взаимно компенсируются. Поэтому полученный вывод о точности нельзя автоматически переносить на произвольные функции и интервалы. В данном задании минимальный шаг $10^{-6}$ выбран не потому, что более крупные шаги нарушили точность, а потому, что именно он является минимальным из предложенного ряда и создаёт достаточную вычислительную нагрузку.

Усреднение трёх запусков снижает влияние фоновых процессов, изменения частоты процессора и работы планировщика, однако полностью не исключает случайные колебания. Оптимальное значение 15 относится к конкретной системе с 16 логическими процессорами и используемой версии .NET. На другом компьютере минимум графика может находиться в иной точке. Следовательно, главным результатом является не универсальное число потоков, а воспроизводимая процедура его определения.

\subsection{Свойства задачи, допускающие многопоточное решение}

Вычисление определённого интеграла методом трапеций хорошо поддаётся распараллеливанию прежде всего благодаря аддитивности:
\[
    \int_a^b f(x)\,dx = \int_a^c f(x)\,dx + \int_c^b f(x)\,dx.
\]
Отрезок можно разделить на любое количество непересекающихся частей, вычислить частичные интегралы независимо и сложить результаты. Рабочим потокам не требуется изменять общую структуру данных или ожидать результаты соседних потоков. Функция и параметры являются общими, но используются только для чтения, а промежуточная сумма хранится локально.

Вторым важным свойством является однородность работы. При постоянном шаге близкие по длине подотрезки содержат приблизительно одинаковое количество трапеций. Это позволяет равномерно распределить нагрузку и уменьшает вероятность того, что основной поток будет долго ожидать один перегруженный участок. Третье свойство --- малый объём объединения результатов: независимо от числа трапеций каждый поток возвращает только одно число. Следовательно, объём синхронизации существенно меньше объёма вычислений.

Вместе с тем многопоточность эффективна не всегда. При крупном шаге вычислений мало, поэтому создание и синхронизация потоков могут занять больше времени, чем сам интеграл. Если вычисляемая функция изменяет общее состояние, требует последовательного доступа к ресурсу или имеет очень различную стоимость на разных участках, независимость частей нарушается либо возникает дисбаланс нагрузки. Порядок сложения чисел с плавающей точкой также может немного изменять итог вследствие конечной точности представления. Поэтому результат следует сравнивать с допустимой погрешностью, а не требовать побитового совпадения.

Таким образом, аддитивность интеграла, независимость подотрезков и компактность частичного результата создают условия для ускорения. Экспериментальные замеры позволяют найти границу, после которой увеличение числа потоков перестаёт быть полезным.
